\enlargethispage{4pt}
\section{Introduction}\label{sec:intro}

Let $X$ be a smooth complex projective variety. The Hodge conjecture asserts
that the cycle class map
\[
  \cl^{p}_{X}\colon \CH^{p}(X)_{\QQ}\longrightarrow \Hdg^{p}(X)
  \coloneqq H^{2p}(X,\QQ)\cap H^{p,p}(X)
\]
is surjective for every $p$. For $p=1$ this is the theorem of Lefschetz on
$(1,1)$ classes.%
\footnote{The rational coefficients are not a convenience. Lefschetz proved the integral statement in degree two, and the integral analogue in higher degree is false: Atiyah and Hirzebruch produced torsion classes of type $(p,p)$ that are not classes of cycles, and Koll\'ar later produced non-torsion ones. So the conjecture has to be stated after tensoring with $\QQ$, and every statement in this paper is a statement about $\QQ$-coefficients.} For $p\ge2$ the conjecture is open, and the two attacks that
are attempted most often are the following.

The first attack is algebraic and runs through abelian varieties. By the
structure theory of Hodge classes on abelian varieties, the Hodge ring of an
abelian variety of Weil type is generated by divisor classes together with
finitely many exceptional classes in middle degree, the Weil classes of
\cite{Weil77}.%
\footnote{Weil introduced them to show that the Hodge ring of an abelian variety is not always generated in degree one, which at the time was an open possibility. He did not conjecture that they are algebraic and did not claim they are not; the examples were offered as a warning, and they have been the standard test case ever since.} Divisor classes are algebraic, so for these varieties the whole
question is whether the Weil classes are algebraic. Markman
\cite{Mar25a,Mar25b} settled this for abelian sixfolds of Weil type of
discriminant $-1$, which are the sixfolds of split Weil type, for every
imaginary quadratic field, and deduced it for all abelian fourfolds of Weil type, for
every discriminant and every field, using a degeneration argument of Schoen
that he quotes; the Hodge conjecture for abelian varieties of dimension at
most five follows. The sheaf he produces has positive rank, is built from the
geometry of the abelian variety itself through a pair of secant sheaves on an
abelian threefold, and satisfies an equivariant semiregularity condition.%
\footnote{The hypothesis he verifies is stronger than injectivity of the
semiregularity map alone, and he asks in \cite[Question 11.4]{Mar25b} whether
it can be relaxed to injectivity on the image of an evaluation map, recording
in his Section 12 that an affirmative answer would carry the method to
abelian varieties of dimension at least eight and to fields of larger degree.
\Cref{cor:kernelfamily} below is a statement about the unrelaxed hypothesis;
\Cref{thm:weakfails} shows that for objects with Chern character concentrated
on the Weil line the relaxed hypothesis fails as well, by the same $n^{2}$
dimensions, and \Cref{prop:evaluation} that for the explicit split object of
\Cref{sec:construction} it fails by the whole kernel; \Cref{thm:linebundle}
that it fails for every object with a line bundle summand off the
polarisation line, at every member and in every dimension, so that no sum of
line bundles can run either argument (\Cref{cor:noline}); and
\Cref{rem:whatmeasures} says what remains, which is the positive rank,
indecomposable case the question was asked for. For that case
\Cref{thm:factor} computes the classes the relaxed hypothesis is about, for
every secant object on an abelian $n$-fold with $n\ge3$, and reduces the
hypothesis for Markman's own candidate in dimension eight to one identity
for a rank one complex on the fourfold; \Cref{thm:candidatefails} shows that
the identity fails, so that candidate does not satisfy the relaxed hypothesis
either.} For
$n\ge4$, and at $n=3$ for the discriminants not covered, the question is
open. It is natural
to hope that the higher cases can be reached by importing an object from an
ambient projective space, restricting it to the dual abelian variety, and
transporting the result to $A$ by the Fourier--Mukai transform. Secant
varieties of the dual are the usual choice of ambient object, because a
dimension count makes the codimension come out right.

The second attack is geometric and runs through induction on dimension. One
places the variety in a family of hyperplane sections, uses the algebraicity
of the Hodge locus \cite{CDK95} to control where the class stays of type
$(p,p)$, appeals to the conjecture in one dimension lower on the sections, and
attempts to transport the resulting cycles back up. Each ingredient here is a
theorem. The question is whether they compose.

This paper proves that neither attack can work, and it is precise about the
reasons, because the reasons are different in kind. The first fails for an
algebraic reason that no refinement of the construction can repair: the
classes the method produces live in the wrong subring, and they live there for
reasons of $\Kx$-equivariance that are visible before any geometry is chosen.
The second fails for a cohomological reason that is elementary once stated:
the restriction map annihilates exactly the classes the induction would need
to carry.

\subsection{The algebraic obstruction}

Throughout, $A$ is a complex abelian variety of dimension $g$, $\Av$ is its
dual, $\cP$ is the Poincar\'e bundle on $A\times\Av$, and
$\Phi_{\cP}\colon D^{b}(\Av)\to D^{b}(A)$ is the Fourier--Mukai transform with
kernel $\cP$. We write $\theta$ for a principal polarisation class on $A$ and
$\theta'$ for the dual class on $\Av$.

\begin{theorem}[Divisor obstruction]\label{thm:divisor}
Let $A$ be a principally polarised complex abelian variety of dimension $g$,
let $i\colon\Av\to\PP^{N-1}$ be any morphism defined by a linear system, and
let $\cF$ be any object of $D^{b}(\PP^{N-1})$. Put
$\cE=\Phi_{\cP}(Li^{*}\cF)$. Then
\[
  \ch_{k}(\cE)\in\QQ\cdot\theta^{k}\qquad\text{for every }k\ge0 .
\]
\end{theorem}

The proof has two ingredients and is short. The Chern character commutes with
derived pullback, which confines $\ch(Li^{*}\cF)$ to the subring of
$H^{*}(\Av,\QQ)$ generated by the hyperplane class, and the hyperplane class of
a linear system on an abelian variety is a multiple of $\theta'$. The
cohomological Fourier--Mukai transform then carries each power of $\theta'$ to
a rational multiple of a power of $\theta$; we prove this in closed form in
\Cref{thm:fmtheta}. Neither ingredient uses anything about $\cF$.

Since a Weil class is by construction not in the subring generated by divisor
classes, the consequence is immediate.

\begin{corollary}\label{cor:neverweil}
With notation as in \Cref{thm:divisor}, suppose in addition that $g=2n$ and
that $A$ is of $(K,-1,n)$-Weil type. Then $\ch_{n}(\cE)$ is a Weil class only
if $\ch_{n}(\cE)=0$.
\end{corollary}

We then give a second proof of \Cref{cor:neverweil} that does not pass through
\Cref{thm:divisor} and that explains, in representation-theoretic terms, why
the first proof had to succeed. The endomorphism $\iota(\tau)$ attached to
$\tau\in\Kx$ acts on $H^{2n}(A,\QQ)$, and the $\Kx$-isotypic decomposition of
that space separates the two objects in question.

\begin{theorem}[Character obstruction]\label{thm:character}
Let $(A,\iota,\eta)$ be of $(K,-1,n)$-Weil type and write
$H^{1}(A,\QQ)\otimes\CC=V_{+}\oplus V_{-}$ for the decomposition into
eigenspaces of $\iota(\sqrt{-d})$. Then
\[
  H^{2n}(A,\QQ)\otimes\CC
  =\bigoplus_{r+s=2n}\textstyle\bigwedge^{r}V_{+}\otimes\bigwedge^{s}V_{-},
\]
with $\iota(\tau)$ acting on the $(r,s)$ summand by
$\tau^{r}\bbar{\tau}^{\,s}$.
The Weil line $\HW(A)$ spans the two extreme summands $(r,s)=(2n,0)$ and
$(0,2n)$, on which $\Kx$ acts by $\tau^{2n}$ and $\bbar{\tau}^{2n}$, while
$\eta^{n}$ lies in the middle summand $(r,s)=(n,n)$, on which $\Kx$ acts by
the norm character $\Nm_{K/\QQ}(\tau)^{n}$. In particular, for every $n\ge1$,
\[
  \HW(A)\cap H^{2n}(A,\QQ)^{\Nm^{n}}=0,
\]
so a class that is a $\Kx$-eigenclass for $\Nm^{n}$ has zero component along
the Weil line.
\end{theorem}

\Cref{thm:character} is what makes the obstruction feel inevitable rather than
accidental. Any construction that is equivariant for the $K$-action, and
whose output is therefore an eigenclass, must choose a character in advance,
and the norm character is the one that the polarisation already occupies. The
Weil line sits at the two ends of the $\Kx$-grading, as far from the middle as
the grading allows. \Cref{fig:charactergrid} draws the grading and marks the
two positions.

It is tempting to look for a third obstruction of a numerical kind, and we
record in \Cref{sec:numerology} why there is none. A secant numerology
producing a class of codimension exactly $n$ on an abelian variety of
dimension $2n$ requires a polarisation on $\Av$ of Euler characteristic
$2n^{2}+2n$, which for $n\ge2$ is smaller than $2^{2n}$ and very much smaller
than $3^{2n}$. Those two quantities are the Euler characteristics at which
the theorem of Lefschetz guarantees that a linear system is free of base
points and that it is very ample, and it is easy to read them as thresholds
below which no such system exists. They are not. Both are sufficient
conditions and neither is necessary: an abelian surface with a polarisation of
type $(1,3)$ has $\chi=3<2^{2}$ and a base-point-free system, and one with a
polarisation of type $(1,5)$ has $\chi=5<3^{2}$ and a very ample system, the
latter being the Horrocks--Mumford embedding in $\PP^{4}$ \cite{HM73,GP98}. We
give the details in \Cref{sec:numerology}, because the erroneous reading
appears in an earlier version of part of this material and we wish to retract
it explicitly. Nothing else in the paper depends on it:
\Cref{thm:divisor,thm:character} are independent of any numerical count.

\subsection{The geometric obstruction}

The second half of the paper concerns induction on dimension. Let
$j\colon Y\hookrightarrow X$ be a smooth ample hypersurface in a smooth
projective variety of dimension $d$, and let $L=c_{1}(\cO_{X}(Y))$.

\begin{theorem}[Hyperplane obstruction]\label{thm:hyperplane}
Let $d=2p$ and let $\alpha\in H^{2p}(X,\QQ)$ be primitive. Then
$j^{*}\alpha=0$ for every smooth ample hypersurface $j\colon Y\hookrightarrow X$.
\end{theorem}

The proof is three lines from the projection formula and the Lefschetz
hyperplane theorem, and we give it in \Cref{sec:hyperplane}. Its force is that
the primitive middle degree is exactly where the Hodge conjecture is hard, and
it is exactly where the restriction map returns nothing. Below the middle
degree the restriction map is injective, but then the return trip is the
problem.

\begin{theorem}[The return trip]\label{thm:returntrip}
Let $2p\le d-2$, so that $j^{*}\colon H^{2p}(X,\QQ)\to H^{2p}(Y,\QQ)$ is an
isomorphism. If $j^{*}\alpha$ is algebraic on $Y$, then $L\alpha$ is algebraic
on $X$. The implication ``$L\alpha$ algebraic $\Rightarrow$ $\alpha$
algebraic'', asked of all smooth projective $X$ and all $\alpha$, is
equivalent to the Lefschetz standard conjecture.
\end{theorem}

Putting the two together gives the statement we want.

\begin{corollary}[No dimensional induction through hyperplane sections]%
\label{cor:noinduction}
Let $X$ be smooth projective of dimension $d$ and let $\alpha\in\Hdg^{p}(X)$.
An argument that deduces the algebraicity of $\alpha$ from the Hodge
conjecture for smooth ample hypersurfaces of $X$, using only $j^{*}$ and
$j_{*}$, returns no information when $2p=d$ and $\alpha$ is primitive, and
otherwise establishes the algebraicity of $L\alpha$ rather than of $\alpha$.
\end{corollary}

\Cref{sec:reduction} then states what is left. By hard Lefschetz every Hodge
class is determined by primitive classes in degrees at most $d$, and the
Lefschetz standard conjecture is the exact obstruction to moving between
degrees algebraically. The conclusion is a clean equivalence,
\Cref{thm:equivalence}: the Hodge conjecture holds for all smooth complex
projective varieties if and only if it holds for primitive classes of middle
degree and the Lefschetz standard conjecture holds. This is a restatement and
not a reduction, and we say so; its value is that it isolates the primitive
middle degree as the only place where new ideas can enter, and
\Cref{thm:hyperplane} says that the classical inductive handle is blind
precisely there.

\subsection{What is not obstructed}\label{ssec:introscope}

We are careful in \Cref{sec:scope} about the limits of these results, because
an obstruction theorem is only useful if its edges are sharp.

\Cref{thm:divisor} does not apply to Markman's construction \cite{Mar25a,Mar25b},
and \Cref{sec:escape} is devoted to saying exactly why, because that section is
the useful one for anybody who wants to build something new. In outline: a
rational class in the Weil line has a $\Kx$-orbit spanning a plane rather than
a line (\Cref{prop:orbitplane}), so any construction reaching it must be
bilinear in a pair of inputs interchanged by $\Gal(K/\QQ)$
(\Cref{cor:orbitcriterion}). Markman's construction is a secant construction,
and it satisfies this because it sites the secant geometry on an abelian
$n$-fold $X$, attaches two conjugate secant sheaves to the two points where a
plane of Hodge classes meets the spinor variety, and pairs them on
$X\times\widehat{X}$. The secant idea is therefore not what fails; siting it on
$\Av$ inside a projective space is what fails, and it fails three times over
(\Cref{prop:whichescaped}). \Cref{q:higher} states what an answer above the
known range would have to supply, and \Cref{sec:construction} supplies three
of its four ingredients in closed form for every $n$ and every $K$: the abelian
$2n$-fold of Weil type is $X\times\widehat{X}$ with a $K$-action written down
by one element of $\NS(X)$ (\Cref{thm:explicitweil}), and the plane is
$P_{t}=\QQ u_{t}\oplus\QQ v_{t}$ with $u_{t}$ and $v_{t}$ the even and odd
parts of $e^{\sqrt{-d}\,t}$ (\Cref{thm:plane}). The third ingredient becomes
an explicit equation: a secant sheaf has rank $a$, $c_{1}=bt$,
$c_{2}=\tfrac{ad+b^{2}}{2}t^{2}$, and Bogomolov discriminant equal to the norm
form $\Nm_{K/\QQ}(b+a\sqrt{-d})$ times $t^{2}$
(\Cref{thm:secantinvariants}), so the numerical obstruction vanishes
identically. At $n=3$ the equation forces the sheaf to be
$\mathcal{I}_{Z}(\Theta)$ with $Z$ a union of $d+1$ translates of the
Abel--Jacobi curve of a genus three Jacobian (\Cref{prop:genus3}), which is
the shape of the object used in \cite{Mar25b}. That equation is then solved in
every dimension. Because the secant plane has no preferred scale, one is free
to pass to a positive multiple of a rational point of it, and after clearing
denominators the Chern characters of the powers of a single line bundle already
span the subring $\QQ[t]$, the matrix that expresses them in the basis
$t^{k}/k!$ being a Vandermonde matrix. A class of positive rank on a smooth
projective variety is the class of a coherent sheaf, so the sheaf exists
(\Cref{lem:posrank} and \Cref{thm:secantexists}), with an explicit integer
witness in line bundles: at $n=4$, $d=3$ and $a=b=1$ the multiple is $6$ and
the witness is
$-23[\cO_{X}]+66[\cO_{X}(t)]-54[\cO_{X}(2t)]+20[\cO_{X}(3t)]-3[\cO_{X}(4t)]$.
The fourth ingredient, the deformation argument, is untouched, and
\Cref{rem:deformation} says so plainly; nothing here is a proof of the Hodge
conjecture at $n\ge4$, and \Cref{rem:thirdclosed} separates what has been
closed from what has not.

\Cref{thm:character} applies to any construction that is $K$-equivariant with
the norm character. A construction that is equivariant for the character
$\tau\mapsto\tau^{2n}$ would land in the right eigenspace, and we record in
\Cref{rem:rightcharacter} what such a construction would have to look like.
We do not know of one, and the difficulty is that $\tau\mapsto\tau^{2n}$ is
not $\QQ$-valued, so it cannot be realised by a correspondence defined over
$\QQ$ without additional structure.

\Cref{thm:hyperplane} concerns hyperplane sections. It does not obstruct
induction through other families, for instance through degenerations, through
blow-ups, or through the Kuga--Satake correspondence, which changes the
variety rather than cutting it. \Cref{sec:ks} sets that construction out in
full and marks its boundary. For a weight-two Hodge structure of K3 type it
transports the question to an abelian variety and turns the class into an
endomorphism, and \Cref{prop:ksnotrestriction} shows why nothing here bears on
it; but on an abelian variety of Weil type the hypothesis $h^{2,0}=1$ holds
only for $n=1$, by \Cref{prop:normpart}, and $n=1$ is exactly the case in
which the Weil line already consists of divisor classes.

\subsection{What the Weil classes look like, and how many there are}

Two further results are proved along the way, and they are the reason the
paper is longer than the first two obstructions require.

\Cref{thm:explicitgens} writes the Weil classes down. On the standard model of
an abelian variety of $(K,-1,n)$-Weil type, built in \Cref{sec:explicit} from
a $K$-vector space of dimension $2n$ and a hermitian form of signature
$(n,n)$, the Weil line has the integral basis
\[
  \omega_{1}=\sum_{|T|\ \mathrm{even}}(-1)^{|T|/2}d^{\,n-|T|/2}\,w^{(T)},
  \qquad
  \omega_{2}=\sum_{|T|\ \mathrm{odd}}(-1)^{(|T|-1)/2}d^{\,n-(|T|+1)/2}\,w^{(T)},
\]
where $T$ runs over the subsets of $\{1,\dots,2n\}$ and $w^{(T)}$ is the
monomial taking $y_{j}$ for $j\in T$ and $x_{j}$ otherwise. Each has $2^{2n-1}$
terms. Written this way the classes are objects one can compute with, and two
things become visible at once: the monomials occurring in them are disjoint
from those occurring in the powers of the polarisation, which gives an
elementary second proof of \Cref{prop:notpoly}; and the balanced signature in
the definition of Weil type is forced, since for any other signature the same
two classes are of type $(p,2n-p)$ and are not Hodge classes at all.

\Cref{thm:hdgcount} counts what is there. For $A$ of $(K,-1,n)$-Weil type with
$\Hg(A)=\SU(V,H)$,
\[
  \dim_{\QQ}\Hdg^{k}(A)=\begin{cases}3,&k=n,\\1,&k\ne n,\end{cases}
  \qquad
  \Hdg^{n}(A)=\QQ\,\eta^{n}\oplus\HW(A),
\]
so the Hodge conjecture for such an $A$, in every degree at once, reduces to
the algebraicity of the single class $\omega_{1}$. That is
\Cref{cor:whatisleft}, and it is what makes the obstruction theorems worth
proving: they are about one explicitly written class and not about a space of
unknown size.

\Cref{thm:divcriterion} then produces that class from divisors whenever a
certain condition holds. Write $R$ for the rational structure on
$\bigwedge^{2}V_{+}\oplus\bigwedge^{2}V_{-}$, the part of $H^{2}(A,\QQ)$ on
which $\Kx$ acts through $\tau^{2}$ and $\bbar{\tau}^{\,2}$. If $R$ contains a
Hodge class whose component in $\bigwedge^{2}V_{-}$ is a nondegenerate form
$\Theta$, then $\HW(A)$ is spanned by the real and imaginary parts of
$\Theta^{n}$, and these are polynomials in two divisor classes, hence
algebraic. On $A=X\times\widehat{X}$ with the $K$-action of
\Cref{sec:construction} the condition holds with
$\Theta=d\beta-\widehat{\beta}-\sqrt{-d}\,\ell$, which is
\Cref{thm:splitclosed}. \Cref{prop:divfails} shows the condition fails
whenever $\Hg(A)$ is the full special unitary group, so the criterion is
exactly a Noether--Lefschetz condition on $R$; and \Cref{prop:splitgeom}
determines what the split members supply: they have trivial discriminant,
whatever polarisation is chosen, they form a locus of codimension
$n(n-1)/2$, and at such a point the three classes
$\beta,\widehat{\beta},\ell$ already exhaust both the Neron--Severi group and
the Hodge classes in the middle degree.

\Cref{lem:Rendomorphism} identifies the criterion: a Hodge class in $R$ is the
same thing as an endomorphism $\mu$ with $\mu\iota(\tau)=\iota(\bbar{\tau})\mu$,
so the criterion asks for a quaternion algebra $(-d,\mu^{2})_{\QQ}$ inside
$\End^{0}(A)$ extending $K$. Those abelian varieties can be written down. For
$B=(-d,b)_{\QQ}$ with $b>0$ acting on $B^{n}$, \Cref{thm:quatweil} shows the
balanced signature is forced rather than imposed, computes the hermitian form
as $\mathrm{diag}(2a_{k}d,-2a_{k}bd)$, identifies the discriminant as the
class of $b^{n}$, and verifies the criterion with $\mu$ left multiplication by
$j$. Since $b$ may be any positive rational, and since discriminants multiply
over products, \Cref{thm:everyfamily} follows: every Weil family, for every
$K$, every $n$ and every discriminant, contains a member whose Weil classes
are algebraic, with a representative that is a polynomial in two divisor
classes. What is then left is propagation and nothing else, which is
\Cref{cor:R2}; and \Cref{cor:exactform} puts that in its sharpest form.
Deligne's theorem makes $\omega_{s}$ absolutely Hodge at every point, so the
Hodge conjecture for abelian varieties of Weil type, for every field, every
dimension and every discriminant, is equivalent to the single implication that
this one absolutely Hodge class is algebraic.

\subsection{What the last implication reduces to}

The set $\Sigma\subseteq\cD_{n,n}$ of points at which $\omega$ is algebraic is
then examined directly, in \Cref{ssec:properness}. Because the relative
Hilbert scheme is proper and the cycle class is locally constant in a flat
family, the part of $\Sigma$ cut out by one tuple of Hilbert polynomials is
Zariski closed; $\Sigma$ is the countable union of those pieces
(\Cref{thm:closedlocus}). Because algebraicity survives an isogeny and
$\SU(V,H)$ acts trivially on the Weil line, $\Sigma$ is stable under
$\SU(V,H)(\QQ)$, and that group has dense orbits, so $\Sigma$ is dense
(\Cref{prop:densealg}). The two facts together give a dichotomy with nothing
between its halves: either $\Sigma$ is everything, and the Hodge conjecture
holds for all abelian varieties of Weil type in the family, or $\Sigma$ has
empty interior and the conjecture fails outside a meagre set
(\Cref{cor:allornothing}).

One further operation preserves $\Sigma$, and it is the only one that moves
between dimensions. If $A_{1}$ and $A_{2}$ are of Weil type for the same $K$,
with parameters $n_{1}$ and $n_{2}$, then $A_{1}\times A_{2}$ is of Weil type
with parameter $n_{1}+n_{2}$, and the Weil class of the product is a fixed
rational combination of external products of the Weil classes of the factors
(\Cref{thm:weilmult}). Written in terms of the $K$-valued invariant
$\omega=\omega_{1}+\sqrt{-d}\,\omega_{2}$ the combination is simply
$\omega(A_{1})\omega(A_{2})=2\,\omega(A_{1}\times A_{2})$, so the Weil class is
multiplicative and algebraicity passes from the factors to the product
(\Cref{cor:prodalg}). At $n=1$ the Weil classes lie in $H^{2}$ and are
therefore algebraic by Lefschetz with no hypothesis at all, so taking $n$
abelian surfaces of Weil type gives, in every dimension and for every $d$, an
abelian $2n$-fold of Weil type whose Weil classes are written as an explicit
integral polynomial in $2n$ divisor classes (\Cref{thm:everyn}). That is the
only statement in this paper which produces an algebraic Weil class in every
dimension while quoting nothing, and its reach is correspondingly small: the
locus of such products has dimension $n$ inside a family of dimension $n^{2}$,
smaller even than the quaternionic locus, and a general member of the family
is simple and lies on none of them (\Cref{rem:prodscope}).

That is what turns the last implication into two statements of finite
character, each about the cycles already constructed. \Cref{thm:degreebound}:
if the Hilbert data of the cycle of \Cref{thm:everyfamily}, transported along
one rational orbit, takes finitely many values, the conjecture follows for the
whole family; at $n=2$, where the quaternionic locus is a divisor, infinitely
many repetitions suffice rather than all but finitely many
(\Cref{rem:ntwo}). \Cref{thm:semiregclosure}: the base point carries a perfect
complex whose Chern character is $N\omega_{1}$ in degree $n$ and zero
elsewhere (\Cref{prop:ktheory}), and if its semiregularity map is injective,
the conjecture again follows for the whole family. Of the two, the second
cannot be run on the object that supplies it: by
\Cref{prop:concentrated} the complex of \Cref{prop:ktheory} has
semiregularity kernel of dimension at least $n^{2}$ at every base point
constructed here, so its hypothesis fails for the object that supplies it
(\Cref{rem:killsreduction}), and by \Cref{cor:kernelfamily} the kernel is a
canonical copy of the tangent space to the family the deformation was to
cross, so no choice of base point inside the family can improve it. The degree bound is untouched, and \Cref{prop:leadingconstant} sharpens
it: the leading coefficient of the Hilbert polynomial is constant along the
orbit, because $G(\QQ)$ preserves both the Weil class and the polarisation,
so what has to be bounded is only the degree of the effective part of the
transported cycle (\Cref{rem:whereunbounded}). The divisorial constructions
cannot supply that bound: products of effective divisors of bounded degree
represent $\omega$ only on a nowhere dense set (\Cref{prop:divroute}), and
\Cref{prop:pencil} shows the degrees growing along an explicit pencil of
quaternionic loci in one fixed polarised lattice. Neither reduction mentions
a general point and neither assumes the variational Hodge conjecture; the
passage from a formal deformation at one point to cycles everywhere, which is
the part that is usually missing, is supplied by \Cref{thm:closedlocus}. We
have not proved either hypothesis, and \Cref{rem:tworemain} says exactly what
each of them asks for.

\Cref{ssec:whatobject} then says what the object of the second criterion has
to be. For a sum of line bundles the semiregularity map is computed in closed
form (\Cref{lem:splitsemireg}): it annihilates every off-diagonal extension
class and the traceless part of every diagonal one, so injectivity forces
multiplicity one, a vanishing of cohomology that the shifts can arrange, and
the injectivity of $(\zeta_{i})\mapsto\sum_{i}e^{\lambda_{i}}\zeta_{i}$.
The Chern character condition then fails outright.
\Cref{thm:positivity}: at any member whose only divisor classes of norm
character are the multiples of $\eta$, if the Chern character of an object is
an effective combination $\sum_{i}m_{i}e^{\lambda_{i}}$ of exponentials of
divisor classes and $\ch_{2}$ is a flat Hodge class, then the coefficient of
$\theta^{+}\theta^{-}$ is $\sum_{i}m_{i}\Nm_{K/\QQ}(u_{i})$ and must vanish;
the norm form is positive definite and the $m_{i}$ are positive, so every
$u_{i}$ is zero and the Weil component of the Chern character is zero as
well. The separation of $\theta^{+}\theta^{-}$ from $\eta^{2}$ that this
needs is \Cref{lem:tensorrank}, and it is a statement about tensor rank: in
the isotypic component for the square of the norm character the first class
has rank one and the second has rank $\binom{2n}{2}$.

The hypothesis on the Chern character holds for every sheaf filtered by line
bundles, for every semi-homogeneous bundle in the sense of Mukai
\cite{Muk78}, and for every iterated extension of those; the hypothesis on
the divisor classes holds at every split member and at every quaternionic base
point. So the candidates one would reach for first do not exist, and the
reason is positivity rather than a computation. What is left must be
indecomposable in the derived category with a class in $K_{0}$ that is not
effective, which is a sharper description of the remaining construction than
we started with.

The same subsection then asks where such an object could live at all, and the
answer is arithmetic. Writing the candidate as a sum of line bundles indexed
by elements $u_{i}$ of the ring of integers of $K$, the Chern character
conditions become $T_{q,r}=\sum_{i}m_{i}u_{i}^{q}\bbar u_{i}^{\,r}=0$ for all
$q,r\le n$ but three, and grouping the $u_{i}$ by norm turns the diagonal part
of that system into $\sum_{j}M_{j}N_{j}^{r}=0$ for $r=1,\dots,n$, where
$N_{1}<\dots<N_{\nu}$ are the norms that occur and $M_{j}$ is the sum of the
multiplicities at level $N_{j}$. Fewer than $n+1$ norms forces every $M_{j}$
to vanish and the rank with it (\Cref{cor:nplusone}); at exactly $n+1$ the
level sums are forced by Lagrange interpolation, and since multiplicity one
allows each element of the ring only once, they are capped by the number of
elements of each norm. \Cref{thm:pterange} runs that competition to the end.
For $4\le n\le8$, six imaginary quadratic fields and norms up to the bounds in
\Cref{sec:verification}, no set of norms meets the two demands, so in that
range the split route closes before any geometry is used. At $n=3$ it does
not: exactly eleven sets of four norms survive, all at $d=1$ or $d=2$, and the
smallest of them is $(9,18,27,36)$ over $\QQ(\sqrt{-2})$, where the level sums
are $(4,-6,4,-1)$ and the object would have rank one. For all eleven the rest
of the system is then shown to have no solution, by an exhaustive search over
sign patterns that covers $4.29\times10^{14}$ cases in the largest instance.

\subsection{Verification}

Every numerical claim in this paper is checked by exact arithmetic, and the
output is reproduced in \Cref{app:scripts}. The scripts verify the closed form
of the cohomological Fourier--Mukai transform on powers of $\theta$, the
stability of $\QQ[\theta]$ under the transform, the numerical question of
\Cref{cor:nonumerical} for $2\le n\le8$, the identity
$\sum_{q}h^{q,q+2}=\binom{4n}{2n-2}$ for the semiregularity target of
\Cref{sec:semireg}, the character separation of \Cref{thm:character}, the
coordinate model and the two generators \eqref{eq:omegagens}, the invariant
count behind \Cref{thm:hdgcount}, the Clifford arithmetic of \Cref{sec:ks},
and the family and secant plane of \Cref{sec:construction}. All arithmetic is
over $\QQ$ or in the exterior algebra of a finite-dimensional vector space
with integer structure constants; no step is numerical in the floating-point
sense. \Cref{sec:lean} describes a second, independent check of the finite
arithmetic by the Lean~4 kernel.

\subsection{Organisation}

\Cref{sec:conventions} fixes conventions. \Cref{sec:weiltype} develops the
structure theory of abelian varieties of Weil type from the definitions,
\Cref{sec:explicit} writes the Weil classes out in coordinates and counts the
Hodge classes, and \Cref{sec:characters} proves \Cref{thm:character}.
\Cref{sec:fm} establishes the closed form of the cohomological Fourier--Mukai
transform, and \Cref{sec:divisor} proves \Cref{thm:divisor} and
\Cref{cor:neverweil}. \Cref{sec:numerology} settles the numerical question in
the negative, and \Cref{sec:semireg} computes the semiregularity target,
explains why it is not the binding constraint either, and then proves
\Cref{thm:rankzero}: no object whose Chern character is concentrated on the
Weil line is semiregular, because the annihilator of the Weil line in
$H^{0,2}$ has dimension $n^{2}$ (\Cref{lem:annihilator}). It then identifies
that annihilator: by \Cref{thm:anntangent} the polarisation carries the
tangent space to the Weil family isomorphically onto it, and by
\Cref{prop:firstorder} a polarisation-preserving deformation kills the whole
Weil line exactly when it commutes with $K$, so the family is also the
first-order Hodge locus of the class. The kernel and the family are therefore
the same space, which is \Cref{cor:kernelfamily}; and \Cref{thm:weakfails}
places that space inside the image of the evaluation map of Hochschild
cohomology, so that the weakened semiregularity criterion of
\cite[Question 11.4]{Mar25b} is unavailable for such objects too; and
\Cref{thm:linebundle} shows, from the infinitesimal rigidity of divisor
classes, that any object with a line bundle summand off the polarisation
line is obstructed along the family in a direction the semiregularity map
cannot see, which removes every sum of line bundles in every dimension. This
is the fourth obstruction, and it is the one that bears directly on the
fourth ingredient of \Cref{rem:deformation}. \Cref{sec:hyperplane}
proves \Cref{thm:hyperplane} and \Cref{thm:returntrip}, and \Cref{sec:ks}
marks the boundary of the one classical route that gets past them.
\Cref{sec:reduction} proves the equivalence, \Cref{sec:escape} identifies which
hypotheses a construction must break, and \Cref{sec:construction} gives a
family of examples in which three of the four ingredients are supplied in
closed form, settles the base case in every family, and closes with
\Cref{ssec:properness} on the structure of the algebraic locus and
\Cref{ssec:whatobject} on the object the second criterion requires, where
\Cref{prop:splitobstruction} computes the first-order obstruction of the
explicit split object against the family: it is nonzero exactly in the
directions normal to the split locus, it lies in the kernel of the
semiregularity map, and it is transverse to the copy of the tangent space
that \Cref{cor:kernelfamily} places there; and \Cref{prop:evaluation} writes
out the evaluation map of Hochschild cohomology on that object and finds that
the whole kernel of the semiregularity map is in its image; and
\Cref{prop:markman8} computes the Chern character of Markman's candidate
object in dimension eight, places it at the point $(1,3)$ of the secant
plane, and reduces its verification to one condition on one sheaf; and
\Cref{thm:factor} moves that condition from the eightfold to the fourfold,
by computing the Hochschild classes that preserve any secant Chern
character, $n^{2}$ of them, so that the weakened criterion for the transform
of $\cF\boxtimes\cF$ becomes the single identity
$\mathrm{At}(\cF)^{2}=-d\Theta^{2}\cdot\id$ in
$\Ext^{2}(\cF,\cF\otimes\Omega^{2})$ for $\cF$ itself
(\Cref{cor:markmancondition}); and \Cref{thm:candidatefails} decides that
identity in the negative, by a local computation at an isolated double point
of the support, so that the candidate of Markman's Section 12 does not
satisfy the weakened criterion, and \Cref{cor:localobstruction} extends the
obstruction to every dimension. What that argument cannot reach is a support
with no such point, and \Cref{ssec:smoothsupport} takes up that case. At a
smooth point the sheaf $\cE xt^{2}$ of the ideal vanishes, so both sides of
the identity are locally zero and the germ comparison has nothing to compare
(\Cref{prop:smoothlocal}); the question moves from a point to the normal
bundle. In exchange the surface itself becomes rigid. Since the Todd class of
an abelian variety is trivial, Grothendieck--Riemann--Roch turns the Chern
character into the whole numerical package of the support, and for
$\ch(\cF)=u_{\Theta}+b\,v_{\Theta}$ on a fourfold it gives
\[
  \chi(\cO_{S})=(b^{2}+d)(3b^{2}-d),\quad
  K_{S}^{2}=3(b^{2}+d)(7b^{2}-d),\quad
  e(S)=3(b^{2}+d)(5b^{2}-3d),
\]
with nothing left free. Because $\Omega^{1}_{S}$ is a quotient of the trivial
bundle $\Omega^{1}_{X}|_{S}$, it is globally generated, so $K_{S}$ is nef and
$e(S)\ge0$; then $K_{S}$ is also big, $S$ is minimal of general type, and the
Bogomolov--Miyaoka--Yau inequality applies. Its defect is
$3e(S)-K_{S}^{2}=24(b^{4}-d^{2})$, so the inequality says exactly $d\le b^{2}$
(\Cref{thm:smoothinvariants}). At the value $b=3$ that the fourfold forces,
this leaves $d\in\{1,3,5,7\}$ (\Cref{cor:fourdiscriminants}): for every
imaginary quadratic field beyond those four there is no secant object of rank
one with smooth support at all, and for those four the invariants are the
six numbers of \Cref{tab:smoothsupport} and nothing else.
The same subsection then identifies what the local argument is really testing.
It is not smoothness and it is not the complete intersection property: at any
point where the support is Cohen--Macaulay the sheaf $\cE xt^{2}$ of its ideal
vanishes, because Auslander--Buchsbaum makes the ideal of projective dimension
one there, so both sides of the identity are locally zero
(\Cref{thm:cmvanishing}). A Macaulay2 computation over ten germs confirms the
line, and shows it is sharp: the fat plane and the cone over the twisted cubic
are not complete intersections and still carry no obstruction, while two
planes meeting at a point, three planes meeting pairwise at a point, and a
plane with an embedded point all do. So the criterion splits cleanly
(\Cref{cor:cmdichotomy}): a support that is not Cohen--Macaulay is removed
outright, and a support that is Cohen--Macaulay has no local obstruction at
all, leaving one global vanishing in
$H^{1}(Z,\cE xt^{1}(\cI_{Z},\cI_{Z}))$. That also explains why Markman's
candidate cannot be repaired by normalising its double points: separating the
branches is what restores the depth, and \Cref{rem:normalisationcm} shows it
moves the Chern character off the secant plane.
A Cohen--Macaulay support of codimension two has an ideal sheaf of projective
dimension one, by the same Auslander--Buchsbaum argument, so it is resolved by
two vector bundles of ranks $r+1$ and $r$. That is a shape one can search.
\Cref{ssec:hilbertburch} does so. When the two bundles are sums of line
bundles the Chern character condition becomes four equations on the twists,
and at $r=1$, the complete intersections, they close outright for every $b$
and every $d$: weighting by $w(x)=x^{2}+d$, which is positive because $d>0$,
turns three of the equations into the statement that two positive measures
share their mass, mean and variance, and a one atom measure has variance zero,
so the two divisor degrees would have to coincide and the resolution would
cancel (\Cref{thm:noci}). At $r=2$ and $r=3$ an exact search over the twists
finds nothing either. At $r=4,5,6$ solutions appear, but only at $d=15$ and
$d=23$ (\Cref{thm:burchsearch}). None of them survives: a morphism from a
secant object that satisfies the criterion to a power of the polarisation is
killed by $H^{2}(\cO)$ (\Cref{lem:annihilation}), and for the cone of a map
between sums of such powers the summand inclusions are not, because the
relevant $H^{2}$ of a nontrivial power vanishes; so no split resolution runs
the criterion at any rank, any twist or any discriminant
(\Cref{thm:nosplit}), and the same holds when the line bundles are replaced
by simple semi-homogeneous bundles, the bundles with scalar Atiyah class
(\Cref{thm:noblocks}). What remains is a support that is either smooth, with
$d$ one of four values, or singular and Cohen--Macaulay, in both cases with a
Hilbert--Burch resolution over a bundle that is not projectively flat and a
restriction $H^{2}(\cO_{X})\to H^{2}(\cO_{Z})$ that is injective
(\Cref{cor:disjointroutes}). \Cref{rem:routemap} writes the two remaining
questions in the form a construction would have to answer.
Before that, \Cref{sec:cmfields} removes the restriction to quadratic
fields. The four theorems that turn the conjecture for a Weil family into a
propagation statement, namely the reduction to one class, the closedness of
the algebraic locus, the density of the rational orbits and the base point,
use nothing about the degree of the field. Carrying them over needs one new
input, a base point, and we supply it for every CM field $F$ of degree $2m$,
every $n$ and every discriminant: the abelian variety is isogenous to a
product of CM abelian $m$-folds, the Weil classes there are polynomials with
rational coefficients in $F$-balanced divisor classes, written out in
\Cref{thm:cmbasepoint}, and they are algebraic by the Lefschetz theorem on
$(1,1)$ classes. The density of the rational orbits needs the hermitian form
to be isotropic over the totally real subfield, which the Hasse principle
gives, and then the unipotent radicals of the parabolic subgroups generate
(\Cref{thm:cmpropagation}). So the Weil classes of the higher CM fields are
not a separate problem standing beside the quadratic ones: they are the same
propagation problem, and the semiregularity criterion, met at one base point
of each family, closes all of them at once. Nothing in that section is conditional on Question 11.4 or on any
secant construction.
\Cref{ssec:closuremap} then sets down, once and in full, what separates all of
this from the conjecture, and \Cref{ssec:closure} decides which parts of it
matter. The implications proved here, together with those the literature
supplies, form a finite rule set, and taking consequences in that rule set is
a computation rather than a judgement. The conjecture is not among the
consequences, which is the machine statement that nothing here proves it.
The route that passes through this paper needs three open statements: a
propagation statement, the one established in \Cref{thm:semiregclosure} and
carried to every CM field in \Cref{thm:cmpropagation}, and by
\Cref{prop:descent} needed only for the split families of the CM fields of
degree at least four, together with the
Hodge classes on abelian varieties that divisor and Weil classes do not
generate, and the conjecture for the varieties that are not abelian
(\Cref{cor:fourremain}). The last of the three is not a remainder: the
varieties that are not abelian include $A\times\PP^{1}$ for every abelian
variety $A$, and the conjecture for $A\times\PP^{1}$ gives it for $A$, so
that statement is equivalent to the whole conjecture (\Cref{prop:f3ishc}).
The statement that belongs in its place is weaker: every Hodge class is
algebraic modulo the images of Hodge classes of abelian varieties under
algebraic correspondences (\Cref{def:f3prime}). It holds on every variety
whose cohomology is reached from abelian varieties in that way, a class
containing curves and abelian varieties and closed under products and
surjective images, in degrees $0$, $2$, $2n-2$ and $2n$, and in dimension at
most three (\Cref{prop:f3prime}); beyond that it is open, and its standard
open case, the square of a K3 surface, turns on the same Kuga--Satake
correspondence as the Mumford classes below (\Cref{rem:f3primeopen}). With
both statements in the rule set there are exactly five minimal sets of open
statements that yield the conjecture (\Cref{thm:closure}). Two are the
conjecture restated: the conjecture for the varieties that are not abelian
alone, and the Lefschetz standard conjecture for every variety together with
the motivatedness of every Hodge class (\Cref{prop:lefmot}). The other three
use the weaker statement, together with the Lefschetz standard conjecture,
with the variational Hodge conjecture for algebraic classes that it implies
(\Cref{prop:lefvhc}, \Cref{thm:vhcab}), or, along the route of this paper,
with the propagation statement and the classes that divisor and Weil classes
do not generate. Only that last route uses a theorem of this paper. Every
statement in the five except the propagation statement is a consequence of
the conjecture, and the Lefschetz standard conjecture, which lies in two of
them and gives the whole abelian half, is the target at the centre
(\Cref{rem:bullseye}). None
of these statements is proved here in general. One new case of the Lefschetz
standard conjecture is: for an abelian scheme over a curve whose invariant
cycles are algebraic, the operator $\Lambda$ is the sum of a relative
Pontryagin product with $\ell^{g-1}/(g-1)!$ and an operator along the base
built from the invariant cycles, so it is algebraic; this covers the total
space of every Mumford family (\Cref{prop:pontryaginlambda},
\Cref{thm:lefschetzfamily}, \Cref{cor:mumfordlefschetz}). The same
decomposition shows that, for such a total space, the Lefschetz standard
conjecture is exactly the statement that algebraicity propagates along the
curve (\Cref{thm:lefschetzpropagation}). With the base points of this paper
that turns each of the two remaining targets into the Lefschetz standard
conjecture for one variety: one Weil family is equivalent to it for the
total space of the family over a curve through a base point
(\Cref{cor:weilequiv}), and the first classes beyond the Weil lines are
equivalent to it for the ninefold $W\times_{C}W$ of a Mumford family
(\Cref{cor:mumfordequiv}). Two statements are proved outright on the way:
the Hodge conjecture for $X_{c}\times X_{c}$ at every CM point $c$ of a Mumford
family, from the Hodge conjecture for abelian fourfolds (\Cref{thm:mumfordcm}),
and the fact that no object meeting the semiregularity criterion is a direct
sum of objects with exterior Ext algebras (\Cref{thm:nosum}). The bounded criterion, which asks that the cycles
representing the class along the orbit have bounded complexity, might look like
a second propagation statement, but by \Cref{thm:p1equivalent} it is
equivalent to the conjecture for the family, a restatement and not a route;
read on the transported cycle it is simply false, because the transport
multiplies the degree of a component by $c^{2n}$ and divides it by the order
of the stabiliser the kernel meets, which is bounded for any component with
finite stabiliser (\Cref{lem:multiplicity}, \Cref{thm:p1forces}). Nor are the
last two statements reductions, to Weil classes and to abelian varieties, in
any sense that would help: the self-product of an abelian fourfold of Mumford
type has eight Hodge classes in $H^{4}$ against six products of divisor
classes, and the two classes left over are Weil classes of nothing
(\Cref{prop:mumfordproduct}); and the $H^{3}$ of a very general quintic
threefold is not of abelian type, because the Hodge cocharacter acts on the
Lie algebra of its Mumford--Tate group with a weight that no tensor
construction on abelian varieties can produce (\Cref{prop:noabeliantype}).
Each of the two is the conjecture for a class of Hodge classes that no Weil
line reaches. The first of them can be pressed one step further. The two
classes on the Mumford fourfold are the real multiplication by a totally real
cubic field on the primitive part of $H^{2}$, and they are algebraic as soon
as the Kuga--Satake class of one K3 surface of Picard number thirteen attached
to the fourfold is; the Clifford relations of the Kuga--Satake construction
are what bring the whole field into play (\Cref{prop:mumfordrm},
\Cref{thm:mumfordks}, \Cref{rem:mumfordks}). That Kuga--Satake class is not known to be algebraic.
The same classes are rigid: in each of the $46$ weight blocks of $HT^{2}$ of
the square the Gram determinant of contraction into a rational exceptional
class is a product of sums of squares and of linear forms that do not vanish
at the conjugates of an irrational element of the cubic field, so the only
first order deformation keeping the class of Hodge type is the direction of
the compact Mumford curve, and no degeneration or larger family reaches it
(\Cref{thm:mumfordrigid}, \Cref{cor:nodegeneration}). Propagation along the
curve follows from one perfect complex at one CM point with $\Ext^{2}$ of
dimension $119$, the least its Chern character allows, or from a degree bound
at infinitely many CM points (\Cref{thm:mumfordnumerical},
\Cref{cor:mumfordbounded}). Such an object is heavily constrained: for a
dense set of classes its self-extensions have dimension at least
$1,16,119,328,560,328,119,16,1$ in degrees $0$ to $8$, the Euler
characteristic forces at least $800$ in odd degrees, its Chern character lies
outside the subring generated by divisor classes, and it extends along no
Hochschild direction but that of the curve (\Cref{thm:mumfordprofile},
\Cref{prop:mumforddivisor}, \Cref{cor:mumfordsplit}). No sequence of cones,
tensor products, pull-backs, push-forwards and Fourier--Mukai functors applied
to line bundles reaches the classes, at any point of the curve, because every
such operation preserves invariance under the Lefschetz group
(\Cref{thm:lefschetzclosure}). Among all complex tori, algebraic or not, the
component of the Hodge locus through the square is the Mumford curve itself,
so no twistor line through a member keeps the classes of Hodge type, and
Markman's twistor mechanism for Weil classes cannot move them; the twistor
lines on which they are of Hodge type carry non-algebraic tori
(\Cref{prop:notwistor}). The square is a holomorphic symplectic variety whose
symplectic class generates a copy of the structure of K3 type in
$H^{1}\otimes H^{1}$, and the classes are exactly the dual forms of that
copy twisted by the real multiplication; so the mechanism of van Geemen and
Rapagnetta \cite{vGR26}, pull-back from a hyperk\"ahler manifold, would need
a rational map onto a symplectic subvariety of a hyperk\"ahler manifold of
dimension at least ten, such as $S^{[n]}$ with $n\ge5$ for a K3 surface of
Picard number thirteen, and no hyperk\"ahler eightfold serves
(\Cref{prop:hksquare}). After Li's proof of the Tate conjecture for
powers of abelian fourfolds over finite fields \cite{Li26}, the classes are
represented by algebraic cycles at every closed point of every reduction of
the curve modulo a prime, and the target is equivalent to a bound on the
Hilbert data of those cycles on a Zariski dense set of closed points of the
arithmetic model, with no lifting step (\Cref{prop:modp},
\Cref{thm:modpdensity}). \Cref{rem:bypassaudit} audits the ten routes that
avoid the Lefschetz standard conjecture, and none is carried out. The five
that are not closed off are not five obligations: three are the target in
other forms and two imply it, so any one would suffice
(\Cref{cor:mumfordone}). And the target is only the first case of the Hodge
classes that divisor and Weil classes do not generate, so a proof of it would
leave open the propagation statement, the other such classes, and the
conjecture for the varieties that are not abelian, which is the whole
conjecture (\Cref{rem:mumfordneeded}, \Cref{prop:f3ishc}).

The propagation statement itself is constrained from several sides, and the
constraints describe the object it asks for closely. Its semiregularity form
forces the evaluation map of Hochschild cohomology to vanish on the
$4n^{2}$-dimensional annihilator of the Weil class, and so forbids the
representative to be a sheaf on any smooth germ of codimension two of its
support (\Cref{thm:p2forces}, \Cref{cor:p2shape}). That annihilator has a
shape: $HH^{1}(A)$ splits as $P\oplus Q$ into the annihilators of the two
conjugate pieces $\alpha_{\pm}$ of the Weil class, each of dimension $2n$,
and the annihilator in degree two is exactly $P\wedge Q$
(\Cref{thm:hhsplitting}). So a complex meeting the criterion has its whole
Hochschild action factoring through
$\bigwedge^{*}P\oplus_{\CC}\bigwedge^{*}Q$, of dimension $2^{2n+1}-1$ in a ring
of dimension $2^{4n}$: it sees the two halves of the class separately and
their interaction not at all. With no hypothesis on semiregularity, every
perfect complex with Chern character on the Weil line has
$\dim\Ext^{2}(E,E)\ge2n(2n-1)$ (\Cref{cor:hhfactor}), and the minimum is
itself enough: a single complex at a single base point with
$\dim\Ext^{2}(E,E)=2n(2n-1)$ has injective semiregularity map and closes the
family (\Cref{thm:p2numerical}). This turns the criterion from a rank into a
number, which is the form in which a candidate can be tested. It fixes more
than itself: by the Hodge--Riemann relations $\chi(E,E)>0$, an indecomposable
summand carries the class, and at $n=2$, and at $n=3$ granting a
compatibility the paper already uses, it has at least three independent
endomorphisms, so no simple object meets the criterion (\Cref{thm:p2endo}).
Finally the support is placed. Two independent jet classes on a complex of finite length
multiply to a nonzero class whenever its Euler characteristic is nonzero,
because on a torus the semiregularity map sees that product through the
Euler characteristic (\Cref{lem:finiteproduct}); applied on a slice
transverse to a component of the support, this shows that every component of
codimension at least two with nonzero multiplicity is a union of cosets of an
abelian subvariety tangent to a $K$-eigenspace, and since the Weil line misses
every class pulled back from such a quotient, the support must have a
component of codimension less than $n$ (\Cref{thm:p2support}). No complex
carried by a cycle that represents the Weil class can serve, at any member of
any family. The hypothesis on the Euler characteristic is sharp: the cone of
the product itself, on the Koszul complex of a point, has Euler characteristic
zero and kills it.

All of this concerns a complex whose Chern character is exactly a multiple of
the Weil class, and no such complex is semiregular (\Cref{thm:p2false}). Its
semiregularity would make it extend along every $K$-linear deformation of the
abelian variety, and those reach complex tori that are not algebraic. On a very
general such Weil torus the only Hodge classes below the top degree are the
Weil classes, there are no subvarieties but points, and every coherent sheaf
has vanishing Chern character in degrees $1$ to $2n-1$: this is Voisin's
theorem in dimension four \cite{Voi02b}, and the same argument through the
Hermite--Einstein metrics of Bando and Siu gives it in every dimension
(\Cref{lem:weiltorushodge}, \Cref{lem:weiltorussubsets},
\Cref{prop:weiltorussheaves}). So the numerical criterion above is never met,
and the results that describe its object are vacuous. What survives is the
corrected statement in which the complex may also carry powers of the
polarisation, which are Hodge on the polarised family and not beyond it; the
propagation theorem holds for it unchanged and it escapes the obstruction;
the semiregular objects in Markman's proofs are not of the pure form either
(\Cref{rem:p2prime}). Its numerical form is computed in
\Cref{thm:p2primenumber}: for a Chern character $N\omega+\sum_{k}c_{k}\theta^{k}$
the annihilator in $HT^{2}$ has dimension $n^{2}(4-\rho)$ for $n\ge3$, with $\rho\le3$ the
rank of the Hankel matrix of the numbers $k!\,c_{k}$, and for a general shape
it is exactly the tangent space of the polarised Weil family; so one complex
with $\dim\Ext^{2}(E,E)=7n^{2}-2n$ at one base point would close the family,
$24$ at $n=2$ and $57$ at $n=3$. Of the shapes with $\rho\le1$, a top power
of $\theta$ and a single exponential $e^{t\theta}$ with $t\theta$ the class of
a line bundle are excluded, by a twist and \Cref{thm:p2false}, and the whole Hochschild profile of
a corrected character is in \Cref{prop:p2primeprofile}. The
Mumford classes escape the obstruction for a different reason: they are
rigid, and their Hodge locus is the Mumford curve.

The secant route, the one this paper spends most of its length narrowing,
lies outside every minimal set: granted every demand it makes, it reaches the
families of trivial discriminant, because the member it is built on is a
product and a product has trivial discriminant, and through descent every
imaginary quadratic family, but no CM field of higher degree
(\Cref{rem:secantplace}). Descent is the one link between the Weil families
(\Cref{prop:descent}): a product with a Weil surface of any discriminant and
a push-forward against its Weil class, which is a divisor class, carry the
Weil classes of one family onto those of every family of lower dimension and
every discriminant, and a scalar extension carries those of a CM field of
degree at least four onto those of the imaginary quadratic fields it
contains. So along the route of this paper the propagation statement is
needed only for the split families of the CM fields of degree at least four,
and the base points of nontrivial discriminant, though they exist, are not
needed on any route. The computation behind the frontier is item
(XXXIII) of \Cref{sec:verification}, repeated in the Lean kernel.
\Cref{sec:final} then states the closure itself, as a single theorem in seven
clauses, every one of them unconditional: the reduction to one class on one
connected domain and its converse, a base point with an explicit cycle in
every family for every field and every discriminant, one criterion for
propagation, which is a number, and one that restates the conjecture, the
fourteen routes that are closed and the structural reason each of them fails,
the same criterion for every CM field, and the three statements that remain
along this route, the last of them the conjecture itself and its weaker form
open, the second not a reduction and its smallest case brought down to a
single Kuga--Satake class, beside the other minimal routes. \Cref{sec:scope} delimits all of it,
\Cref{sec:verification} describes
the verification and \Cref{sec:lean} the Lean certificate.
\Cref{app:exterior} carries the coordinate computation behind the transform,
\Cref{app:cdk} the theorem of Cattani, Deligne and Kaplan on the locus of
Hodge classes, \Cref{app:albert} the place of Weil type inside Albert's
classification, and \Cref{app:scripts} the output of the scripts.
